% Part: first-order-logic
% Chapter: syntax-and-semantics
% Section: formation-sequences

\documentclass[../../../include/open-logic-section]{subfiles}

\begin{document}

\olfileid[fr]{pl}{syn}{fseq}

\olsection{Suites de formation}

Définir les !!{formula}s inductivement et démontrer leurs propriétés
par induction constituent deux démarches complémentaires, élégantes
et efficaces. Il peut cependant être utile de suivre leur construction
pas à pas, en partant des éléments les plus simples. C'est ce que
permet la notion de \emph{suite de formation}.

\begin{defn}[Suites de formation des formules]
\ollabel{defn:fseq-frm}
Une suite finie $\tuple{!A_0,\dotsc,!A_n}$ de chaînes de symboles
du langage~$\Lang L_0$ est une \emph{suite de formation} de $!A$ si
$!A \ident !A_n$ et si, pour tout $i \leq n$, ou bien $!A_i$ est
une formule atomique, ou bien il existe $j,k < i$ tels que l'un
des cas suivants se présente :
\begin{enumerate}
    \tagitem{prvNot}{$!A_i \ident \lnot !A_j$.}{}%
    \tagitem{prvAnd}{$!A_i \ident (!A_j \land !A_k)$.}{}%
    \tagitem{prvOr}{$!A_i \ident (!A_j \lor !A_k)$.}{}%
    \tagitem{prvIf}{$!A_i \ident (!A_j \lif !A_k)$.}{}%
    \tagitem{prvIff}{$!A_i \ident (!A_j \liff !A_k)$.}{}%
\end{enumerate}
\end{defn}

\begin{ex}
\[
\tuple{
    \Obj p_0,
    \Obj p_1,
    (\Obj p_1 \land \Obj p_0),
    \lnot (\Obj p_1 \land \Obj p_0)
}
\]
est une suite de formation de
$\lnot (\Obj p_1 \land \Obj p_0)$, tout comme
\[
\tuple{
    \Obj p_0,
    \Obj p_1,
    \Obj p_0,
    (\Obj p_1 \land \Obj p_0),
    (\Obj p_0 \lif \Obj p_1),
    \lnot (\Obj p_1 \land \Obj p_0)
}.
\]
%
Comme le montre le second exemple, une suite de formation peut
contenir des éléments superflus : des formules redondantes ou
qui ne contribuent pas à la construction.
\end{ex}

\begin{prop}
\ollabel{prop:formed}
Toute !!{formula}~$!A$ de~$\Frm[L_0]$ admet une suite de formation.
\end{prop}

\begin{proof}
Supposons $!A$ atomique. Alors la suite~$\tuple{!A}$ est une suite
de formation de~$!A$.
%
Supposons maintenant que $!B$ et~$!C$ admettent respectivement les
suites de formation $\tuple{!B_0,\dotsc,!B_n}$ et
$\tuple{!C_0,\dotsc,!C_m}$.
%
\begin{enumerate}
  \tagitem{prvNot}{Si $!A \ident \lnot !B$,
      alors $\tuple{!B_0,\dotsc,!B_n,\lnot !B_n}$
      est une suite de formation de~$!A$.}{}
  \tagitem{prvAnd}{Si $!A \ident (!B \land !C)$,
      alors $\tuple{!B_0,\dotsc,!B_n,!C_0,\dotsc,!C_m,(!B_n \land !C_m)}$
      est une suite de formation de~$!A$.}{}
  \tagitem{prvOr}{Si $!A \ident (!B \lor !C)$,
      alors $\tuple{!B_0,\dotsc,!B_n,!C_0,\dotsc,!C_m,(!B_n \lor !C_m)}$
      est une suite de formation de~$!A$.}{}
  \tagitem{prvIf}{Si $!A \ident (!B \lif !C)$,
      alors $\tuple{!B_0,\dotsc,!B_n,!C_0,\dotsc,!C_m,(!B_n \lif !C_m)}$
      est une suite de formation de~$!A$.}{}
  \tagitem{prvIff}{Si $!A \ident (!B \liff !C)$,
      alors $\tuple{!B_0,\dotsc,!B_n,!C_0,\dotsc,!C_m,(!B_n \liff !C_m)}$
      est une suite de formation de~$!A$.}{}
\end{enumerate}
Le principe d'induction sur les !!{formula}s permet de conclure que
toute !!{formula} admet une suite de formation.
\end{proof}

On peut également démontrer la réciproque. Ce résultat est important :
il montre que les deux définitions des formules sont équivalentes,
c'est-à-dire qu'elles donnent le même ensemble. On peut donc aussi
démontrer des théorèmes sur les formules par récurrence ordinaire
sur la longueur des suites de formation.

\begin{lem}
\ollabel{lem:fseq-init}
Supposons que $\tuple{!A_0,\dotsc,!A_n}$ soit une suite de formation
de~$!A_n$ et que $k \leq n$. Alors $\tuple{!A_0,\dotsc,!A_k}$ est
une suite de formation de~$!A_k$.
\end{lem}

\begin{proof}
La démonstration est laissée en exercice.
\end{proof}

\begin{thm}
\ollabel{thm:fseq-frm-equiv}
$\Frm[L_0]$ est l'ensemble de toutes les chaînes de symboles du
langage~$\Lang L_0$ qui admettent une suite de formation.
\end{thm}

\begin{proof}
Soit $F$ l'ensemble des chaînes de symboles du langage~$\Lang L_0$
qui admettent une suite de formation. La
\olref[pl][syn][fseq]{prop:formed} donne $\Frm[L_0] \subseteq F$ ;
il reste à démontrer l'inclusion réciproque.

Supposons que $!A$ admette une suite de formation
$\tuple{!A_0,\dotsc,!A_n}$. Montrons que $!A \in \Frm[L_0]$ par
récurrence forte sur le dernier indice~$n$. L'hypothèse de récurrence
est la suivante : toute chaîne de symboles admettant une suite de
formation $\tuple{!B_0,\dotsc,!B_m}$ avec $m < n$ appartient à
$\Frm[L_0]$. Par définition, $!A_n$ est atomique, ou bien il existe
$j,k < n$ tels que l'un des cas suivants se présente :
\begin{enumerate}
    \tagitem{prvNot}{$!A_n \ident \lnot !A_j$.}{}%
    \tagitem{prvAnd}{$!A_n \ident (!A_j \land !A_k)$.}{}%
    \tagitem{prvOr}{$!A_n \ident (!A_j \lor !A_k)$.}{}%
    \tagitem{prvIf}{$!A_n \ident (!A_j \lif !A_k)$.}{}%
    \tagitem{prvIff}{$!A_n \ident (!A_j \liff !A_k)$.}{}%
\end{enumerate}
Raisonnons par cas. Si $!A_n$ est atomique, alors
$!A_n \in \Frm[L_0]$. Supposons plutôt que
$!A \ident (!A_j \land !A_k)$. D'après le
\olref[pl][syn][fseq]{lem:fseq-init},
$\tuple{!A_0,\dotsc,!A_j}$ et $\tuple{!A_0,\dotsc,!A_k}$
sont des suites de formation de $!A_j$ et de~$!A_k$ respectivement.
Ce sont des préfixes propres de la suite de formation de~$!A$ :
leurs derniers indices $j$ et $k$ sont strictement inférieurs
à~$n$. L'hypothèse de récurrence donne donc
$!A_j,!A_k \in \Frm[L_0]$. Par définition d'une !!{formula},
$(!A_j \land !A_k)$ appartient lui aussi à $\Frm[L_0]$.
Les autres cas se traitent de la même manière.
\end{proof}
\begin{editorial}
La suite indexée de $0$ à $n$ a $n+1$ termes. La source confond
par endroits longueur et dernier indice. La preuve ci-dessus
utilise uniformément le dernier indice pour la récurrence.
Le symbole isolé $\equiv$ de la source y est également remplacé
par le symbole d'identité syntaxique $\ident$.
\end{editorial}

\end{document}
